\section{Análisis exploratorio}
\label{sec:eda}

El análisis exploratorio se realizó sobre las \(1000\) observaciones del
dataset antes de ajustar modelos predictivos. Su finalidad fue reconocer
la estructura de la cartera, identificar distribuciones asimétricas,
examinar posibles diferencias entre buenos y malos créditos y formular
hipótesis que posteriormente pudieran evaluarse mediante inferencia
estadística y regresión.

En esta etapa, los resultados se interpretan como asociaciones
marginales. Una diferencia observada entre grupos no implica que la
variable mantenga el mismo efecto cuando se consideran simultáneamente
otras características del solicitante y de la operación. Tampoco permite
establecer relaciones causales.

Las visualizaciones que respaldan esta lectura se presentan en la
Sección~\ref{sec:visualizaciones}. La estadística descriptiva completa,
incluidos los estadísticos de posición, dispersión, asimetría y curtosis,
se desarrolla en la Sección~\ref{sec:descriptiva}.

\subsection{Distribución de la variable respuesta}

La muestra contiene \(700\) operaciones clasificadas como buen riesgo y
\(300\) como mal riesgo. En términos relativos, la tasa observada de mal
crédito es:

\[
\widehat{p}
=
\frac{300}{1000}
=
0.300.
\]

Por tanto, la clase mayoritaria representa el \(70\,\%\) de la muestra y
la clase de interés el \(30\,\%\). Se trata de un desbalance moderado,
pero suficiente para que la exactitud global pueda ofrecer una impresión
incompleta del desempeño predictivo.

Por ejemplo, una regla trivial que clasificara todos los registros como
buen riesgo alcanzaría una exactitud del \(70\,\%\), aunque no
identificaría ningún caso de mal crédito. Esta observación justifica que
la evaluación posterior incluya \textit{recall}, precisión, matrices de
confusión, ROC--AUC y costos de clasificación, en lugar de utilizar
únicamente el \textit{accuracy}.

\begin{hallazgo}
La clase positiva no es mayoritaria. En consecuencia, el desempeño del
modelo debe evaluarse por su capacidad para identificar malos créditos y
por el costo de los errores, no solo por el porcentaje total de
clasificaciones correctas.
\end{hallazgo}

\subsection{Propósito del préstamo}

La tasa global de mal crédito no se distribuye de manera uniforme entre
los propósitos declarados. La categoría Educación presenta la mayor tasa
entre los grupos con tamaño moderado, con \(22\) malos créditos de
\(50\) observaciones, equivalente al \(44.0\,\%\). Le siguen Auto nuevo,
con \(38.0\,\%\), y Negocio, con \(35.1\,\%\).

En el extremo inferior aparecen Radio/TV, con una tasa de
\(22.1\,\%\), Auto usado, con \(16.5\,\%\), y Reentrenamiento, con
\(11.1\,\%\). La diferencia entre Educación y la tasa global es de
\(14.0\) puntos porcentuales, mientras que Auto usado se sitúa
\(13.5\) puntos porcentuales por debajo del promedio de la cartera.

Sin embargo, la comparación requiere considerar el tamaño de cada
categoría. Otros, Reparaciones, Electrodomésticos y Reentrenamiento
contienen menos de treinta observaciones. En estos grupos, una variación
de pocos casos produce cambios importantes en la tasa calculada.

La Tabla~\ref{tab:eda-proposito} presenta los resultados exploratorios
principales.

\begin{table}[H]
    \centering
    \caption{Tasa observada de mal crédito según el propósito del
    préstamo.}
    \label{tab:eda-proposito}

    \begin{tabularx}{0.92\textwidth}{
        @{}
        X
        >{\centering\arraybackslash}p{1.5cm}
        >{\centering\arraybackslash}p{1.5cm}
        >{\centering\arraybackslash}p{2.4cm}
        >{\centering\arraybackslash}p{2.5cm}
        @{}
    }
        \toprule
        Propósito
        & Total
        & Malos
        & Tasa de mal crédito
        & Diferencia frente al global \\
        \midrule

        Educación
        & 50
        & 22
        & \(44.0\,\%\)
        & \(+14.0\) pp \\

        Otros
        & 12
        & 5
        & \(41.7\,\%\)
        & \(+11.7\) pp \\

        Auto nuevo
        & 234
        & 89
        & \(38.0\,\%\)
        & \(+8.0\) pp \\

        Reparaciones
        & 22
        & 8
        & \(36.4\,\%\)
        & \(+6.4\) pp \\

        Negocio
        & 97
        & 34
        & \(35.1\,\%\)
        & \(+5.1\) pp \\

        Electrodomésticos
        & 12
        & 4
        & \(33.3\,\%\)
        & \(+3.3\) pp \\

        Mobiliario/equipos
        & 181
        & 58
        & \(32.0\,\%\)
        & \(+2.0\) pp \\

        Radio/TV
        & 280
        & 62
        & \(22.1\,\%\)
        & \(-7.9\) pp \\

        Auto usado
        & 103
        & 17
        & \(16.5\,\%\)
        & \(-13.5\) pp \\

        Reentrenamiento
        & 9
        & 1
        & \(11.1\,\%\)
        & \(-18.9\) pp \\

        \bottomrule
    \end{tabularx}

    \begin{minipage}{0.92\textwidth}
        \footnotesize
        \textit{Nota.} La tasa global de mal crédito es
        \(30.0\,\%\). Las categorías Otros, Reparaciones,
        Electrodomésticos y Reentrenamiento tienen menos de
        \(30\) observaciones; sus tasas deben interpretarse con mayor
        cautela.
    \end{minipage}
\end{table}

Estas diferencias sugieren una asociación entre el propósito del
préstamo y la clasificación de riesgo. No obstante, la evidencia
exploratoria no permite concluir todavía que la asociación sea
estadísticamente significativa. La prueba correspondiente se presenta
en la Sección~\ref{sec:inferencia}.

\subsection{Duración del crédito}

La duración muestra uno de los contrastes exploratorios más claros entre
las dos clases. En los buenos créditos, la mediana es de \(18\) meses y
la media de \(19.21\) meses. En los malos créditos, la mediana aumenta a
\(24\) meses y la media a \(24.86\) meses.

También cambia la dispersión. El rango intercuartílico es de \(12\)
meses entre los buenos créditos y de \(24\) meses entre los malos. Por
tanto, el grupo de mal riesgo no solo presenta contratos más prolongados,
sino también una mayor heterogeneidad en sus plazos.

Desde una perspectiva de riesgo, una duración mayor amplía el horizonte
durante el cual pueden cambiar los ingresos del solicitante, sus gastos,
el entorno económico o su capacidad de pago. Sin embargo, esta lectura
es una interpretación financiera plausible y no una relación causal
demostrada por los datos.

\subsection{Monto del crédito}

El monto también presenta valores superiores entre los casos de mal
riesgo. La mediana es de \(2244\) DM para los buenos créditos y de
\(2574.5\) DM para los malos. Las medias son \(2985.46\) DM y
\(3938.13\) DM, respectivamente.

La diferencia entre medias es mayor que la diferencia entre medianas,
debido a la fuerte asimetría de la distribución y a la presencia de
operaciones de importe elevado. Además, el rango intercuartílico del
monto es de \(2259.25\) DM para los buenos créditos y de \(3789\) DM
para los malos.

Estas características indican que la comparación no debe descansar
únicamente en la media. Por ello, el análisis posterior utiliza la
mediana, métodos no paramétricos y la transformación
\variable{log_monto}, según el objetivo de cada procedimiento.

Existe, además, un solapamiento considerable entre ambas clases. Hay
buenos créditos de monto elevado y malos créditos de monto reducido. En
consecuencia, el monto aporta información sobre el perfil de riesgo, pero
no permite separar por sí solo las dos clases.

\subsection{Edad del solicitante}

Los casos clasificados como mal riesgo corresponden, en promedio, a
solicitantes algo más jóvenes. La mediana de edad es de \(31\) años en
el grupo de mal riesgo y de \(34\) años en el grupo de buen riesgo. Las
medias respectivas son \(33.96\) y \(36.22\) años.

La diferencia es menor que la observada en la duración y existe un
amplio solapamiento entre ambas distribuciones. La edad se considera,
por tanto, una señal potencialmente relevante, pero insuficiente para
sostener una conclusión aislada. Su contribución deberá evaluarse dentro
del modelo multivariable.

Además, por tratarse de una característica personal, cualquier uso
operacional de la edad requeriría revisar su pertinencia regulatoria,
ética y de equidad.

\subsection{Carga de la cuota}

La variable \variable{tasa_cuota} representa una escala ordinal de carga
financiera, no el porcentaje exacto de la cuota sobre el ingreso. El
nivel \(4\), correspondiente a la categoría superior de la escala,
concentra el \(53.0\,\%\) de los malos créditos y el \(45.3\,\%\) de
los buenos.

Esta comparación no significa que la tasa de mal crédito del nivel
\(4\) sea del \(53.0\,\%\). El porcentaje está calculado dentro de cada
clase y muestra que la categoría superior tiene mayor presencia relativa
entre los casos de mal riesgo.

El patrón es compatible con la hipótesis de que una mayor carga de pago
puede reducir el margen financiero disponible del solicitante. No
obstante, la interpretación debe mantenerse en términos ordinales,
porque el dataset no proporciona el porcentaje exacto de ingreso
comprometido.

\subsection{Relaciones monotónicas entre variables}

La matriz de correlaciones de Spearman permite evaluar asociaciones
monotónicas entre las variables numéricas, las escalas ordinales y la
respuesta binaria.

La relación más alta aparece entre duración y monto:

\[
\rho_{\text{duración,monto}}=0.62.
\]

Esto indica que los créditos de mayor importe tienden a contratarse por
plazos más extensos. La relación es suficientemente relevante para
justificar una revisión posterior de multicolinealidad, aunque no implica
que ambos predictores sean equivalentes.

En relación con la respuesta:

\[
\rho_{\text{duración,mal crédito}}=0.21,
\]

\[
\rho_{\text{monto,mal crédito}}=0.09,
\]

\[
\rho_{\text{edad,mal crédito}}=-0.11.
\]

La duración presenta la asociación monotónica marginal más clara con el
mal crédito. El monto tiene una asociación positiva mucho menor, mientras
que la edad muestra una asociación negativa débil. Las demás variables
numéricas y ordinales presentan coeficientes próximos a cero.

Estos resultados no deben utilizarse para descartar automáticamente
predictores. Una correlación marginal reducida puede coexistir con una
asociación condicional relevante cuando se controlan otras variables,
especialmente en presencia de categorías, interacciones o relaciones no
lineales.

\subsection{Síntesis exploratoria}

La Tabla~\ref{tab:sintesis-eda} resume los patrones que orientan las
etapas siguientes.

\begin{table}[H]
    \centering
    \caption{Síntesis de los principales patrones exploratorios.}
    \label{tab:sintesis-eda}

    \begin{tabularx}{0.94\textwidth}{
        @{}
        >{\raggedright\arraybackslash}p{2.8cm}
        X
        >{\raggedright\arraybackslash}p{4.1cm}
        @{}
    }
        \toprule
        Dimensión
        & Evidencia exploratoria
        & Implicancia analítica \\
        \midrule

        Clase de riesgo
        & \(70\,\%\) buenos y \(30\,\%\) malos.
        & No evaluar el modelo únicamente mediante exactitud. \\

        Propósito
        & Tasas entre \(11.1\,\%\) y \(44.0\,\%\), con tamaños de
          grupo heterogéneos.
        & Aplicar prueba de independencia y considerar la estabilidad
          de las categorías pequeñas. \\

        Duración
        & Mediana de \(18\) meses en buenos y \(24\) meses en malos.
        & Comparar grupos y evaluar su asociación condicional. \\

        Monto
        & Valores mayores y más dispersos en los malos créditos.
        & Utilizar medidas robustas, escala logarítmica y contraste no
          paramétrico. \\

        Edad
        & Mediana de \(34\) años en buenos y \(31\) en malos.
        & Evaluar si la asociación persiste dentro del modelo
          multivariable. \\

        Nivel de cuota
        & Mayor presencia del nivel \(4\) entre los malos créditos.
        & Interpretar como escala ordinal y no como porcentaje exacto. \\

        Duración--monto
        & Correlación de Spearman de \(0.62\).
        & Revisar redundancia y multicolinealidad. \\

        \bottomrule
    \end{tabularx}
\end{table}

En conjunto, el EDA señala que la duración, el monto, la edad, la carga
de la cuota y el propósito contienen información potencialmente útil
para caracterizar el riesgo crediticio. Sin embargo, ninguna de estas
variables separa de forma completa las clases y todas las comparaciones
siguen siendo marginales.

Las diferencias observadas se someterán a contrastes frecuentistas y a
modelos multivariables. Esta transición es necesaria para determinar qué
patrones se sostienen después de cuantificar la incertidumbre y controlar
simultáneamente por otras características.
% Contenido previsto:
% - Balance de la clase.
% - Distribución por propósito.
% - Duración, monto, edad y variables de capacidad de pago.
% - Patrones marginales, valores extremos y asimetría.
% - Hipótesis exploratorias, sin inferencia causal.
